\section{Livrable 3 -- Modes de marches et d'arrêts}

\begin{UPSTIinfor}{Rappel : GMMA}
Le Grafcet de Modes de Marches et d'Arrêts (GMMA) permet de structurer les
différents modes de fonctionnement d'un système automatisé : production
normale, marches de préparation/vérification, marches de fin de cycle,
arrêts d'urgence, mise en/hors énergie, etc. Chaque mode se caractérise par
des conditions d'entrée, des conditions de sortie, et les actions qu'il
autorise ou impose.
\end{UPSTIinfor}

\begin{UPSTIactivite}[][Modes de fonctionnement de la cellule]
    \UPSTIquestion{Identifier les différents \textbf{modes de
    fonctionnement} nécessaires pour exploiter la cellule robotique en
    tenant compte du besoin exprimé (traitement des gobelets dans les 4
    cuves) et des exigences habituelles de sécurité et d'exploitation d'une
    installation robotisée.}

    \UPSTIquestion{Pour chaque mode retenu, préciser :
    \begin{itemize}
        \item les \textbf{conditions d'entrée} dans le mode ;
        \item les \textbf{conditions de sortie} du mode ;
        \item les \textbf{actions} réalisées par l'automate, le robot et le
        convoyeur pendant ce mode.
    \end{itemize}}

    \UPSTIquestion{Représenter l'articulation de ces modes sous la forme
    d'un \textbf{GMMA} (ou d'un schéma équivalent), en identifiant clairement
    les transitions entre modes.}
\end{UPSTIactivite}
